\section*{Chapitre \ref{ch:b1:p-block-13-15} --- Le bloc p, I~: les groupes du bore, du carbone et de l'azote}

\begin{solution}{exo:b1:p-block-13-15:1}
\ce{NH4NO3}~: $-$III dans \ce{NH4+}, $+$V dans \ce{NO3-}. \ce{N2O5}~: $+$V.
\ce{HNO2}~: $+$III. \ce{N2H4}~: $-$II. \ce{Mg3N2}~: $-$III.
\end{solution}

\begin{solution}{exo:b1:p-block-13-15:2}
\ce{BF3}~: bore à trois liaisons et six électrons, trigonal plan.
\ce{NH3}~: trois liaisons et un \omterm{def:b1:lewis-resonance-vsepr:lewis}{doublet non liant}. Le \omterm{def:b1:lewis-resonance-vsepr:lewis}{doublet non liant}
de N forme une \omterm{def:b1:lewis-resonance-vsepr:lewis-acid}{liaison dative} avec B~: \ce{BF3} est l'\omterm{def:b1:lewis-resonance-vsepr:lewis-acid}{acide de Lewis},
\ce{NH3} la \omterm{def:b1:lewis-resonance-vsepr:lewis-acid}{base de Lewis}~; dans l'adduit, les deux atomes sont
tétraédriques.
\end{solution}

\begin{solution}{exo:b1:p-block-13-15:3}
Avec la soude~: \ce{CO2 + 2NaOH -> Na2CO3 + H2O},
\ce{SiO2 + 2NaOH -> Na2SiO3 + H2O} et
\ce{P4O10 + 12NaOH -> 4Na3PO4 + 6H2O}. L'alumine réagit avec un acide,
\ce{Al2O3 + 6HCl -> 2AlCl3 + 3H2O}, et avec une base,
\ce{Al2O3 + 2NaOH + 3H2O -> 2NaAl(OH)4}.
\end{solution}

\begin{solution}{exo:b1:p-block-13-15:4}
$\Delta_r G^\circ = 2 \times (-16.45) = \qty{-32.9}{kJ/mol}$, $\log K =
32.9/5.708 = 5.76$, soit $K = 10^{5.76}$.
\end{solution}

\begin{solution}{exo:b1:p-block-13-15:5}
Cycle de six motifs \ce{SiO3^2-}~: \ce{Si6O18^12-}. Chaîne double, pour
quatre siliciums~: deux à deux sommets partagés (3 O, charge $-2$
chacun), deux à trois ($2.5$ O, charge $-1$ chacun)~: \ce{Si4O11^6-}.
\end{solution}

\begin{solution}{exo:b1:p-block-13-15:6}
\qty{550}{kg} de \ce{Al(OH)3} représentent $550/78.0 = \qty{7.05}{kmol}$,
qui donnent
\qty{3.53}{kmol} de \ce{Al2O3}, soit \qty{360}{kg}~; à \qty{92}{\percent}~:
\qty{331}{kg}.
\end{solution}

\begin{solution}{exo:b1:p-block-13-15:7}
Le bore possède une orbitale vacante et accepte le \omterm{def:b1:lewis-resonance-vsepr:lewis}{doublet non liant}
d'un ion hydroxyde pris à l'eau~: \ce{B(OH)3 + 2H2O <=> B(OH)4- + H3O+}.
Le proton libéré provient d'une molécule d'eau, et non de l'acide
borique.
\end{solution}

\begin{solution}{exo:b1:p-block-13-15:8}
N passe de $-$III à $+$II, soit cinq électrons par \ce{NH3}~; chaque
\ce{O2} en capte
quatre. $4 \times 5 = 5 \times 4$~: \ce{4NH3 + 5O2 -> 4NO + 6H2O}.
\end{solution}

\begin{solution}{exo:b1:p-block-13-15:9}
\ce{NH3 + HNO3 -> NH4NO3}. Une tonne représente $10^6/80.0 = \qty{1.25e4}{mol}$~;
\qty{1.25e4}{mol} d'ammoniac pour neutraliser et \qty{1.25e4}{mol} pour
fabriquer l'acide nitrique~: $2.50 \times 10^4 \times 17.0 = \qty{425}{kg}$.
\end{solution}

\begin{solution}{exo:b1:p-block-13-15:10}
Les atomes de la deuxième \omterm{def:b1:periodicity:table}{période} sont petits~: leurs orbitales
\textit{p} se recouvrent bien latéralement et forment des liaisons $\pi$
fortes, si bien que \ce{N#N} et \ce{O=C=O} sont de petites molécules
stables. Les atomes de la troisième \omterm{def:b1:periodicity:table}{période} sont plus gros~; leur
recouvrement $\pi$ est médiocre, et ils préfèrent former davantage de
liaisons simples~: tétraèdres \ce{P4}, réseau de \ce{SiO4}.
\end{solution}

\begin{solution}{exo:b1:p-block-13-15:11}
$E^\circ(\ce{PbO2}/\ce{Pb^2+}) = 1.46 > E^\circ(\ce{Cl2}/\ce{Cl-}) = 1.36$~V~:
\ce{PbO2 + 4H+ + 2Cl- -> Pb^2+ + Cl2 + 2H2O}. Le plomb(IV) est instable
par rapport au plomb(II) (paire inerte)~; le silicium(II) n'est pas un
état stable et
\ce{SiO2} n'est pas un \omterm{def:b1:nernst:couple}{oxydant}.
\end{solution}

\begin{solution}{exo:b1:p-block-13-15:12}
$160 \times 17.0/14.0 = \qty{194}{Mt}$ d'ammoniac, contenant $194 \times
3.0/17.0 = \qty{34}{Mt}$ d'hydrogène, produit surtout à partir du gaz
naturel (méthane et vapeur d'eau).
\end{solution}

\begin{solution}{pb:b1:p-block-13-15:1}
\textbf{1.} :N\,$\equiv$\,N: --- une triple liaison, très forte, sans
polarité.
\textbf{2.} 0, $-$III, $+$II, $+$IV, $+$V, et $-$III/$+$V dans \ce{NH4NO3}.
\textbf{3.} Les bactéries des racines des légumineuses, la foudre, le
procédé Haber--Bosch.
\textbf{4.} Rompre la triple liaison exige un \omterm{def:b1:elementary-steps:catalyst}{catalyseur} que les plantes
ne possèdent pas~; les bactéries dotées de l'enzyme nitrogénase en
disposent.
\textbf{5.} L'azote est réduit (de 0 à $-$III), l'hydrogène oxydé (de 0
à $+$I).
\textbf{6.} \ce{N2 + 3H2 <=> 2NH3}.
\textbf{7.} $K = 10^{2 \times 16.45/5.708} = 10^{5.76}$.
\textbf{8.} À \qty{25}{\celsius}, la réaction est beaucoup trop lente~;
le \omterm{def:b1:elementary-steps:catalyst}{catalyseur} n'agit qu'à chaud, et les conditions sont un compromis
discuté dans le volume de deuxième année.
\textbf{9.} \qty{58.8}{kmol} d'ammoniac~: \qty{29.41}{kmol} de \ce{N2},
soit \qty{824}{kg}~; \qty{88.2}{kmol} de \ce{H2}, soit \qty{176}{kg}.
\textbf{10.} $V(\ce{N2}) = 29.4 \times 10^3 \times 8.314 \times 298.15/10^5 =
\qty{729}{m^3}$~; air~: $729/0.78 = \qty{935}{m^3}$.
\textbf{11.} Du gaz naturel (vaporeformage du méthane), avec du \ce{CO2}
comme sous-produit.
\textbf{12.} Injecté dans le sol, c'est un engrais bon marché, mais c'est
un gaz toxique~; les solides (urée, nitrate d'ammonium) sont plus faciles
à stocker et à épandre.
\textbf{13.} N~: $-$III $\to$ $+$II (5 électrons), O~: 0 $\to$ $-$II (4 par
\ce{O2})~: \ce{4NH3 + 5O2 -> 4NO + 6H2O}.
\textbf{14.} \ce{2NO + O2 -> 2NO2}.
\textbf{15.} \ce{3NO2 + H2O -> 2HNO3 + NO}~: une \omterm{def:b1:e-ph-diagrams:disproportionation}{dismutation} (N passe de
$+$IV à
$+$V et $+$II).
\textbf{16.} \ce{NH3 + 2O2 -> HNO3 + H2O}.
\textbf{17.} $2 \times 58.8 \times 32.0 = \qty{3.76}{t}$.
\textbf{18.} Elle rend l'oxydation rapide et sélective en faveur de NO,
au lieu du \ce{N2}, plus stable.
\textbf{19.} \ce{NH3 + HNO3 -> NH4NO3}.
\textbf{20.} La moitié.
\textbf{21.} \qty{29.4}{kmol} de \ce{NH4NO3}, soit \qty{2.35}{t}.
\textbf{22.} $28.0/80.0 = \qty{35}{\percent}$.
\textbf{23.} Il contient un \omterm{def:b1:nernst:couple}{oxydant} (le nitrate) et un \omterm{def:b1:nernst:couple}{réducteur}
(l'ammonium) dans le même \omterm{def:b1:crystals-metals:crystal}{cristal}~: chauffé ou mêlé à des combustibles,
il peut se décomposer de façon explosive.
\textbf{24.} $58.8 \times 63.0 =$ \textbf{\qty{3.7}{t} d'acide nitrique
par tonne d'ammoniac}.
\end{solution}
